\hypertarget{microservices-and-grpc-in-python}{%
\chapter{Microservices and gRPC in
Python}\label{microservices-and-grpc-in-python}}

\hypertarget{why-grpc}{%
\subsection{79.1 Why gRPC?}\label{why-grpc}}

gRPC is a high-performance RPC framework developed by Google. *
\textbf{Protocol Buffers}: A binary serialization format that is much
faster than JSON. * \textbf{HTTP/2}: Supports multiplexing and
server-side streaming.

\hypertarget{implementing-grpc-in-python}{%
\subsection{79.2 Implementing gRPC in
Python}\label{implementing-grpc-in-python}}

We use the \passthrough{\lstinline!grpcio!} and
\passthrough{\lstinline!protobuf!} libraries to generate C++ accelerated
Python code from \passthrough{\lstinline!.proto!} definitions. This
allows for near-zero-copy communication between microservices written in
different languages.

\begin{center}\rule{0.5\linewidth}{0.5pt}\end{center}

\hypertarget{phase-xix-scientific-computing-internals}{%
\section{Phase XIX: Scientific Computing
Internals}\label{phase-xix-scientific-computing-internals}}

Python's dominance in science is due to its ability to wrap
high-performance Fortran and C libraries.

\hypertarget{chapter-86-numpy-internals-memory-strides-and-ufuncs}{%
\chapter{Chapter 86: NumPy Internals: Memory Strides and
UFuncs}\label{chapter-86-numpy-internals-memory-strides-and-ufuncs}}

\hypertarget{the-ndarray-c-struct}{%
\subsection{\texorpdfstring{86.1 The \texttt{ndarray}
C-Struct}{86.1 The ndarray C-Struct}}\label{the-ndarray-c-struct}}

A NumPy array is a C-struct that points to a block of data. *
\textbf{Data}: Pointer to the raw memory. * \textbf{Dimensions}: Shape
of the array. * \textbf{Strides}: The number of bytes to skip in memory
to reach the next element in a dimension. This allows for \(O(1)\)
reshaping and slicing without copying data.

\hypertarget{universal-functions-ufuncs}{%
\subsection{86.2 Universal Functions
(UFuncs)}\label{universal-functions-ufuncs}}

UFuncs are C-loops that operate on \passthrough{\lstinline!ndarray!}
data. They handle type dispatching and SIMD acceleration (Chapter 71)
automatically.

\begin{center}\rule{0.5\linewidth}{0.5pt}\end{center}
